\clearpage
\section{\no Bounded Synthesis Slices: a Calculus for \texorpdfstring{$\mathit{MinSynSlice}$}{MinSynSlice}}
\label{app:bounded-syn-slices}
\mechfile{Slicing/Synthesis/\\BoundedSynthesis,\\SynSliceCalc}

\subsection*{Biased precision}

Read $\Phi \preceq_a \gamma$ as: \emph{wherever $\Phi$ pins a concrete
type at an assumption, $\gamma$ must record at least as much there};
positions where $\Phi$ is $\gap$ impose no constraint, and $\gamma$ is
free to be $\gap$ at any position regardless of $\Phi$. So $\Phi$ acts
as a partial floor on assumption usage: a lower bound at the positions
$\Phi$ chooses to constrain, silent elsewhere.

\smallskip

Formally:
\[
  \inference{}{[] \preceq_a []}
  \qquad
  \inference{\tau \neq \gap & \tau \sqsubseteq \tau' & \Gamma \preceq_a \Gamma'}
    {(\tau,\, \Gamma) \preceq_a (\tau',\, \Gamma')}
\]
\[
  \inference{\Gamma \preceq_a \Gamma'}{(\gap,\, \Gamma) \preceq_a (\tau',\, \Gamma')}
  \qquad
  \inference{\Gamma \preceq_a \Gamma'}{(\tau,\, \Gamma) \preceq_a (\gap,\, \Gamma')}
\]
The asymmetry between the two $\gap$ rules is why $\preceq_a$ is not a
preorder: both $\gap \preceq_a \tau$ and $\tau \preceq_a \gap$ hold,
though $\gap \neq \tau$.

\subsection*{Judgements}

\fbox{$D \blacktriangleleft \upsilon \;\Rightarrow\; \psi \dashv \gamma$}\ \ \ Fixed-assumption slice: for a synthesis derivation $D : \synthesis{e}{\tau}$ and query $\upsilon \sqsubseteq \tau$, realised type $\psi$ (with $\upsilon \sqsubseteq \psi \sqsubseteq \tau$) and used-assumption slice $\gamma \sqsubseteq \Gamma$.

\medskip

\fbox{$\Phi \vdash D \blacktriangleleft \upsilon \;\Rightarrow\; \psi \dashv \gamma$}\ \ \ Bounded slice: as above, additionally subject to the biased-precision restriction $\Phi \preceq_a \gamma$, where the input context $\Phi$ is a lower bound on the extracted assumption-slice $\gamma$.

\medskip

Setting $\Phi = \bot$ (all-$\gap$) recovers an unconditional
minimality constraint, so the bounded judgement at $\bot$
characterises $\mathit{MinSynSlice}$ directly.

\subsection*{Rules}

The two cases that exhibit the essential pattern are the variable rule (where $\Phi$ is consulted to inflate the recorded assumption slice) and the pair rule (where $\Phi$ is propagated across the two sub-derivations to defeat overlap). The hole and unit leaves are also shown, as they illustrate the trivial case where $\Phi$ is simply absorbed.

\paragraph{Hole and unit.}
\[
  \inference[$\mathit{Hole}$]
    {}
    {\hole[] \blacktriangleleft \gap \;\Rightarrow\; \gap \dashv \bot_\Gamma}
  \qquad
  \inference[$\mathit{Hole}_\Phi$]
    {}
    {\Phi \vdash \hole[] \blacktriangleleft \gap \;\Rightarrow\; \gap \dashv \Phi}
\]
\[
  \inference[$\mathit{Unit}$]
    {}
    {\ast \blacktriangleleft \ast \;\Rightarrow\; \ast \dashv \bot_\Gamma}
  \qquad
  \inference[$\mathit{Unit}_\Phi$]
    {}
    {\Phi \vdash \ast \blacktriangleleft \ast \;\Rightarrow\; \ast \dashv \Phi}
\]

\paragraph{Variable.}
\[
  \inference[$\mathit{Var}$]
    {x{:}\tau \in \Gamma & \upsilon \neq \gap}
    {x \blacktriangleleft \upsilon \;\Rightarrow\; \upsilon \dashv \bot_\Gamma[x \mapsto \upsilon]}
  \qquad
  \inference[$\mathit{Var}_\Phi$]
    {x{:}\tau \in \Gamma & \upsilon \neq \gap}
    {\Phi \vdash x \blacktriangleleft \upsilon \;\Rightarrow\; \Phi(x) \sqcup \upsilon \dashv \Phi[x \mapsto \Phi(x) \sqcup \upsilon]}
\]

\paragraph{Pair.}
\[
  \inference[$\&$]
    {D_1 \blacktriangleleft \upsilon_1 \;\Rightarrow\; \psi_1 \dashv \gamma_1 &
     D_2 \blacktriangleleft \upsilon_2 \;\Rightarrow\; \psi_2 \dashv \gamma_2}
    {(D_1,\,D_2) \blacktriangleleft \upsilon_1 \times \upsilon_2 \;\Rightarrow\; \psi_1 \times \psi_2 \dashv \gamma_1 \sqcup \gamma_2}
\]
\[
  \inference[$\&_\Phi$]
    {(\Phi \sqcup \gamma_2) \vdash D_1 \blacktriangleleft \upsilon_1 \;\Rightarrow\; \psi_1 \dashv \gamma_1 &
     (\Phi \sqcup \gamma_1) \vdash D_2 \blacktriangleleft \upsilon_2 \;\Rightarrow\; \psi_2 \dashv \gamma_2}
    {\Phi \vdash (D_1,\,D_2) \blacktriangleleft \upsilon_1 \times \upsilon_2 \;\Rightarrow\; \psi_1 \times \psi_2 \dashv \gamma_1 \sqcup \gamma_2}
\]

\subsection*{Ruling out the product counterexample}

Counterexample~\ref{lem:ce-prod-not-closed-min} of Section~\ref{sec:ce-prod-not-closed-min} showed that pairing independently-chosen minimal slices does not yield a minimal product: pairing an $\mathbf{if}\,\mathbf{true}\,\mathbf{then}\,x\,\mathbf{else}\,y$ slice with a slice of $x$ aliases the assumption $x$, after which the inner if-slice is no longer minimal in the joined context. The unbounded rules $\&$ and $\mathit{Var}$ have no way to see this aliasing: each sub-derivation is sliced in isolation, then the assumption slices are joined post-hoc.

The bounded rule $\&_\Phi$ side-steps this. When slicing $D_1$, the lower bound is set to $\Phi \sqcup \gamma_2$, which already \textit{contains} the second component's contribution; via $\mathit{Var}_\Phi$, every use of an aliased assumption $x$ now records $\Phi(x) \sqcup \upsilon$, so $D_1$ is sliced knowing the full assumption available to it, not just its own demand. The same applies symmetrically to $D_2$. Sliced under these inflated bounds, neither sub-derivation can produce a $\gamma_i$ whose join with the other's $\gamma_j$ admits further shrinkage at an aliased position -- the joint minimum is reached directly.

%%% Local Variables:
%%% mode: LaTeX
%%% TeX-master: "PolymorphicTypeSlicing"
%%% End:
